% Part: turing-machines
% Chapter: machines-computations
% Section: halting-states

\documentclass[../../../include/open-logic-section]{subfiles}

\begin{document}

\olfileid[tr]{tur}{mac}{hal}
\olsection{Durma Durumları}

\begin{explain}
Makinelerimizi yalnızca yerine getirilecek bir yönerge kalmadığında duracak
biçimde tanımlamış olsak da Turing makinelerinin yaygın gösterimlerinde
$h\in Q$ olan özel bir \emph{durma durumu}~$h$ bulunur.

Durma durumunun ardındaki fikir basittir: Makine işlemini bitirdiğinde
(girdiyi kabul etmeye hazır olduğunda ya da çıktıyı yazmayı bitirdiğinde),
durduğu bir $h$ durumuna geçer. Bazı makinelerde biri girdiyi kabul eden,
öteki girdiyi reddeden iki durma durumu vardır.
\end{explain}

\begin{ex}\emph{Durma Durumları}.
Bu kavramı açıklamak için çift sayı makinesinin değiştirilmiş bir biçimiyle
başlayalım. Girdi çift olduğunda makinenin $q_0$ durumunda durması yerine,
makineyi bir durma durumuna gönderen bir yönerge ekleyebiliriz.
\[
\begin{tikzpicture}[->,>=stealth',shorten >=1pt,auto,node distance=2.8cm,
                    semithick]
  \tikzstyle{every state}=[fill=none,draw=black,text=black]

  \node[initial, state]         (A)                     {$q_0$};
  \node[state]         (B) [right of=A] {$q_1$};
  \node[state]         (C) [below of=A] {$h$};

  \path (A) edge [bend left] node {\TMtrans{\TMstroke}{\TMstroke}{R}} (B)
  	    edge node {\TMtrans{\TMblank}{\TMblank}{N}} (C)
        (B) edge [loop above] node {\TMtrans{\TMblank}{\TMblank}{R}} (B)
            edge [bend left] node {\TMtrans{\TMstroke}{\TMstroke}{R}} (A);
\end{tikzpicture}
\]

Örneği biraz daha geliştirelim. Makine girdinin tek olduğunu saptadığında
hiç durmaz. Döngü yönergesini, bir ret durumu~$r$'ye gitme yönergesiyle
değiştirerek makineye bir \emph{ret} durumu ekleyebiliriz.
\[
\begin{tikzpicture}[->,>=stealth',shorten >=1pt,auto,node distance=2.8cm,
                    semithick]
  \tikzstyle{every state}=[fill=none,draw=black,text=black]

  \node[initial,state]         (A)                     {$q_0$};
  \node[state]         (B) [right of=A] {$q_1$};
  \node[state]         (C) [below of=A] {$h$};
  \node[state]         (D) [below of=B] {$r$};

  \path (A) edge [bend left] node {\TMtrans{\TMstroke}{\TMstroke}{R}} (B)
  	    edge node {\TMtrans{\TMblank}{\TMblank}{N}} (C)
        (B) edge node {\TMtrans{\TMblank}{\TMblank}{N}} (D)
            edge [bend left] node {\TMtrans{\TMstroke}{\TMstroke}{R}} (A);
\end{tikzpicture}
\]
\end{ex}

\begin{explain}
Özel bir durma durumu eklemek, makinenin belirli girdileri ne zaman kabul ya
da reddettiğini açıkça gösterdiği bunun gibi durumlarda yararlı olabilir.
Ancak özel bir durma durumu eklemenin hesaplama gücünü artırmadığını belirtmek
önemlidir. Benzer biçimde, daha az biçimsel bir durma kavramının da kendine
özgü üstünlükleri vardır. Bu bölümde şimdiye dek kullanılan durma tanımı,
\emph{Durma Problemi} kanıtını sezgisel ve gösterimi kolay kılar. Bu nedenle
özgün tanımımızı kullanmayı sürdürüyoruz.
\end{explain}

\end{document}
